% 根本を疑え
% Reflections on Trusting Trust
% Ken Thompson
\section{Bayesian OODA}
\subsection{Bayesian Vision}
\subsection{OODA}
%OODAループの提唱者であるジョン・ボイド（John Boyd）大佐は、ループの4要素（Observe, Orient, Decide, Act）の中で「Orient（情勢判断／適応・方向づけ）」を中核（重力中心）と位置づけました。単なる「状況の分析」ではなく、入力された観測データを解釈・統合し、行動や意思決定を方向づける認知フィルターです。
%
%ボイドが遺したスケッチにおいて、Orientは相互に作用し合う**5つの要素**で構成されています。
%
%---
%
%### Orient を構成する5つの要素
%
%```text
%  ┌────────────────────────────────────────────────────────┐
%  │                         Orient                         │
%  │                                                        │
%  │     遺伝的特質 (Genetic Heritage)                      │
%  │                     │                                  │
%  │     文化的伝統 (Cultural Traditions)                   │
%  │                     │                                  │
%  │  ── 新しい情報 ─── 過去の経験 (Previous Experience) ── │ ──> Decide / Act
%  │    (New Info)       │                                  │
%  │                   分析と統合 (Analysis & Synthesis)    │
%  │                                                        │
%  └────────────────────────────────────────────────────────┘
%
%```
%
%#### 1. 遺伝的特質（Genetic Heritage）
%
%* 生物学的な制約、本能的反射、知覚の物理的限界、個人の生得的な素質や認知特性。
%* 人間が感覚器を通して受け取れる情報帯域や、恐怖・闘争・逃走などの生理的反応パターンを規定します。
%
%#### 2. 文化的伝統（Cultural Traditions）
%
%* 所属する国家、組織、共同体、時代背景から刷り込まれた価値観、規範、暗黙の前提。
%* 「何が正しく、何が危険で、何が受け入れ可能か」という思考のバイアスや共通の文脈を形成します。
%
%#### 3. 過去の経験（Previous Experience）
%
%* 個人の成功・失敗体験、受けた訓練、蓄積されたエピソード記憶。
%* 目の前の事象に対して「以前のあれと同じパターンだ」と直感的に照合（メンタルシミュレーション）する際の参照フレームとなります。
%
%#### 4. 新しい情報（New Information）
%
%* 「Observe（観察）」のフェーズから絶え間なく流入してくる最新のシグナルや環境変化。
%* 既存のメンタルモデルを揺さぶり、破壊または再編を促すトリガーとして機能します。
%
%#### 5. 分析と統合（Analysis & Synthesis）
%
%* **破壊（分析: Analysis）**: 既存の概念や固定観念を要素ごとに一度分解すること。
%* **創造（統合: Synthesis）**: 分解された要素と流入した新しい情報を組み合わせ、現在の現実に即したまったく新しい概念（メンタルモデル）を組み上げること。
%* ボイドの論文『破壊と創造（Destruction and Creation）』に基づく核心プロセスであり、環境に適応し続けるための認知的再構築エンジンです。
%
%---
%
%### Orient が持つ本質的役割
%
%* **決定をバイパスする即応性（暗黙の誘導）:**
%Orientが極限まで洗練されると、「Decide（明示的な意思決定）」を経由せず、Orientから直接「Act（行動）」へと暗黙的なコマンドが送られる（Implicit Guidance & Control）ようになります。これが達人の「直感的な即応行動」の正体です。
%* **認知の歪み（バイアス）の源泉であり、その打破装置:**
%遺伝・文化・経験は強固なフィルターとして働き、容易に認知バイアスを生みます。ボイドは、絶え間ない「新しい情報」の取り込みと「分析と統合（破壊と創造）」のサイクルによって、その歪みを自覚的に更新し続けることこそが環境適応の唯一の道であると説きました。
\subsection{Double Feedback Loop Learning}
